% ==============================================================================
% 1. KEPENDUDUKAN DAN DEMOGRAFI
% ==============================================================================

\section{Kependudukan dan Demografi}

Publikasi ini menyajikan gambaran komprehensif mengenai kondisi demografi, sosial, ekonomi, dan fasilitas di Desa Popontolen berdasarkan kompilasi data potensi dari 7 (tujuh) wilayah Jaga untuk tahun 2024 dan 2025. Data ini diharapkan dapat menjadi rujukan utama bagi Pemerintah Desa dalam merumuskan perencanaan pembangunan desa yang berbasis data akurat (\textit{evidence-based policy}).

\subsection{Penduduk Menurut Wilayah Jaga dan Jenis Kelamin}

Pada tahun 2024, total penduduk Desa Popontolen tercatat sebanyak \textbf{1.621 jiwa} (845 laki-laki dan 776 perempuan), dan pada tahun 2025 tercatat sebanyak \textbf{1.570 jiwa} yang terdiri dari 816 laki-laki (52,0\%) dan 754 perempuan (48,0\%). Rasio jenis kelamin (\textit{Sex Ratio}) berada pada angka \textbf{108,2}. Secara spasial, \textbf{Jaga 7} menampung 488 jiwa atau sekitar \textbf{31,1\%} dari total populasi desa.

\vspace{0.3cm}

\begin{table}[!htbp]
\centering
\caption{Jumlah Penduduk Menurut Wilayah Jaga dan Jenis Kelamin (2024--2025)}
\label{tab:penduduk_jaga}
\small
\rowcolors{3}{bpslight}{white}
\begin{tabularx}{\textwidth}{@{} l >{\raggedleft\arraybackslash}X >{\raggedleft\arraybackslash}X >{\raggedleft\arraybackslash}X >{\raggedleft\arraybackslash}X >{\raggedleft\arraybackslash}X >{\raggedleft\arraybackslash}X @{}}
\toprule
\rowcolor{bpsnavy}
\textcolor{white}{\textbf{Wilayah Jaga}} & 
\textcolor{white}{\textbf{Laki (2024)}} & 
\textcolor{white}{\textbf{Perem (2024)}} & 
\textcolor{white}{\textbf{Total 2024}} & 
\textcolor{white}{\textbf{Laki (2025)}} & 
\textcolor{white}{\textbf{Perem (2025)}} & 
\textcolor{white}{\textbf{Total 2025}} \\
\midrule
Jaga 1 & 75 & 77 & \textbf{152} & 76 & 78 & \textbf{154} \\
Jaga 2 & 103 & 88 & \textbf{191} & 103 & 88 & \textbf{191} \\
Jaga 3 & 73 & 80 & \textbf{153} & 70 & 80 & \textbf{150} \\
Jaga 4 & 119 & 118 & \textbf{237} & 105 & 102 & \textbf{207} \\
Jaga 5 & 104 & 82 & \textbf{186} & 98 & 81 & \textbf{179} \\
Jaga 6 & 111 & 101 & \textbf{212} & 105 & 96 & \textbf{201} \\
Jaga 7 & 260 & 230 & \textbf{490} & 259 & 229 & \textbf{488} \\
\midrule
\rowcolor{bpsborder!40}
\textbf{TOTAL DESA} & \textbf{845} & \textbf{776} & \textbf{1.621} & \textbf{816} & \textbf{754} & \textbf{1.570} \\
\bottomrule
\end{tabularx}
\end{table}

\subsection{Penduduk Menurut Kelompok Umur}

Struktur umur penduduk Desa Popontolen didominasi oleh kelompok usia produktif (15--59 tahun) yang mencapai \textbf{942 jiwa (64,4\%)} pada tahun 2025. Kelompok usia anak-anak (0--14 tahun) tercatat sebanyak 309 jiwa (21,1\%), dan penduduk lanjut usia (60 tahun ke atas) sebanyak 212 jiwa (14,5\%). Rasio Ketergantungan (\textit{Dependency Ratio}) Desa Popontolen berada pada \textbf{55,3\%}.

\vspace{0.3cm}

\begin{table}[!htbp]
\centering
\caption{Jumlah Penduduk Menurut Kelompok Umur (2024--2025)}
\label{tab:kelompok_umur}
\small
\rowcolors{3}{bpslight}{white}
\begin{tabularx}{\textwidth}{@{} l >{\raggedleft\arraybackslash}X >{\raggedleft\arraybackslash}X >{\raggedleft\arraybackslash}X >{\raggedleft\arraybackslash}X >{\raggedleft\arraybackslash}X >{\raggedleft\arraybackslash}X @{}}
\toprule
\rowcolor{bpsnavy}
\textcolor{white}{\textbf{Wilayah}} & 
\multicolumn{2}{c}{\textcolor{white}{\textbf{Usia 0--14 Thn}}} & 
\multicolumn{2}{c}{\textcolor{white}{\textbf{Usia 15--59 Thn}}} & 
\multicolumn{2}{c}{\textcolor{white}{\textbf{Usia 60+ Thn}}} \\
\rowcolor{bpsnavy}
\textcolor{white}{\textbf{Jaga}} & \textcolor{white}{\textbf{2024}} & \textcolor{white}{\textbf{2025}} & \textcolor{white}{\textbf{2024}} & \textcolor{white}{\textbf{2025}} & \textcolor{white}{\textbf{2024}} & \textcolor{white}{\textbf{2025}} \\
\midrule
Jaga 1 & 21 & 22 & 109 & 110 & 22 & 22 \\
Jaga 2 & 38 & 38 & 19 & 19 & 25 & 25 \\
Jaga 3 & 31 & 32 & 86 & 87 & 32 & 32 \\
Jaga 4 & 54 & 40 & 152 & 138 & 31 & 29 \\
Jaga 5 & 35 & 30 & 136 & 129 & 15 & 20 \\
Jaga 6 & 39 & 39 & 120 & 119 & 53 & 44 \\
Jaga 7 & 102 & 108 & 350 & 340 & 38 & 40 \\
\midrule
\rowcolor{bpsborder!40}
\textbf{TOTAL DESA} & \textbf{320} & \textbf{309} & \textbf{972} & \textbf{942} & \textbf{216} & \textbf{212} \\
\bottomrule
\end{tabularx}
\end{table}

\subsection{Profil Inklusivitas: Penyandang Disabilitas}

Pendataan potensi desa tahun 2025 mencatat keberadaan \textbf{23 orang penyandang disabilitas} di Desa Popontolen yang tersebar di seluruh wilayah Jaga. Ragam disabilitas yang terdata meliputi: Disabilitas Mental (9 jiwa / 39,1\%), Disabilitas Fisik (7 jiwa / 30,4\%), Sensorik Netra (3 jiwa / 13,0\%), Sensorik Wicara (3 jiwa / 13,0\%), dan Sensorik Rungu (1 jiwa / 4,3\%).

\vspace{0.3cm}

\begin{table}[!htbp]
\centering
\caption{Sebaran Penyandang Disabilitas Menurut Ragam Disabilitas per Jaga (2025)}
\label{tab:disabilitas}
\small
\rowcolors{3}{bpslight}{white}
\begin{tabularx}{\textwidth}{@{} l >{\raggedleft\arraybackslash}X >{\raggedleft\arraybackslash}X >{\raggedleft\arraybackslash}X >{\raggedleft\arraybackslash}X >{\raggedleft\arraybackslash}X >{\raggedleft\arraybackslash}X >{\raggedleft\arraybackslash}X @{}}
\toprule
\rowcolor{bpsnavy}
\textcolor{white}{\textbf{Wilayah Jaga}} & 
\textcolor{white}{\textbf{Fisik}} & 
\textcolor{white}{\textbf{Mental}} & 
\textcolor{white}{\textbf{Intelek.}} & 
\textcolor{white}{\textbf{Netra}} & 
\textcolor{white}{\textbf{Wicara}} & 
\textcolor{white}{\textbf{Rungu}} & 
\textcolor{white}{\textbf{Total}} \\
\midrule
Jaga 1 & 1 & 0 & 0 & 0 & 0 & 0 & \textbf{1} \\
Jaga 2 & 2 & 1 & 0 & 1 & 0 & 0 & \textbf{4} \\
Jaga 3 & 1 & 1 & 0 & 0 & 0 & 0 & \textbf{2} \\
Jaga 4 & 0 & 2 & 0 & 0 & 1 & 0 & \textbf{3} \\
Jaga 5 & 0 & 1 & 0 & 0 & 2 & 1 & \textbf{4} \\
Jaga 6 & 2 & 0 & 0 & 0 & 0 & 0 & \textbf{2} \\
Jaga 7 & 1 & 4 & 0 & 2 & 0 & 0 & \textbf{7} \\
\midrule
\rowcolor{bpsborder!40}
\textbf{TOTAL DESA} & \textbf{7} & \textbf{9} & \textbf{0} & \textbf{3} & \textbf{3} & \textbf{1} & \textbf{23} \\
\bottomrule
\end{tabularx}
\end{table}
